\section{\texorpdfstring{\(L0\)}{L0}: expression universality}\label{sec:L0}

This section fixes the single primitive and the precise boundary of the imported
expression theorem.  Throughout, \(\operatorname{Log}\) is the principal complex
logarithm on
\[
  \Omega_{\mathrm{Log}}:=\CC\setminus(-\infty,0],
  \qquad -\pi<\operatorname{Im}\operatorname{Log}z<\pi,
  \qquad e^{\operatorname{Log}z}=z .
\]

\begin{definition}[EML operator]\label{def:eml}
The Exp-Minus-Log operation is the partial binary operation
\[
  \eml\colon\CC\times\Omega_{\mathrm{Log}}\to\CC,
  \qquad \eml(x,y)=e^x-\operatorname{Log}y .
\]
An EML term is interpreted only at those inputs for which every second argument
of every occurrence of \(\eml\) lies in \(\Omega_{\mathrm{Log}}\).  In particular,
nonzero points on the chosen branch cut are outside the complex domain unless a
real-domain calculation separately keeps the relevant second arguments in the
positive real axis.
\end{definition}

\begin{lemma}[Real-domain safety for the internal EML core]
\label{lem:internal-core-branch-domain-safety}
Every occurrence of \(\operatorname{Log}\) in the internally proved real EML core
of \Cref{lem:eml-printed-core,lem:eml-real-witness-library} is evaluated at a
positive real number on the stated real domain of that macro.  Consequently the
internal core identities are identities in the exact partial semantics of
\Cref{def:eml}; they do not use any value of the principal logarithm on the
negative real branch cut.
\end{lemma}

\begin{proof}
For \(E(u)=\eml(u,1)\), the only logarithm input is \(1>0\).  For
\(L(u)=\eml(1,E(\eml(1,u)))\), the domain condition is \(u>0\), so the inner
second input \(u\) is positive and the outer second input
\(E(\eml(1,u))=e^{e-\log u}=e^e/u\) is positive.  In
\(D(a,b)=\eml(L(a),E(b))\), the displayed domain has \(a>0\), and the second
input of the outer \(\eml\) is \(E(b)=e^b>0\).  The macros
\(P,N,A,S\) in \Cref{lem:eml-printed-core} are compositions of these three
forms; their stated domains include exactly the positivity hypotheses just
listed, and \(P(u)=e^u-u\ge1>0\) supplies the additional positive input used by
\(N\) and \(A\).

The witnesses in \Cref{lem:eml-real-witness-library} are then built from
\(E,L,D,P,N,A,S\) together with the auxiliary positive inputs explicitly stated
there: \(E(v)=e^v>0\), positive rational constants, products of positive terms in
the positive-domain multiplier, nonzero denominators only through positive
squares, and the final square-root argument restricted to \((0,\infty)\).  Hence
every logarithm input appearing in those calculations is a positive real number,
which lies in \(\Omega_{\mathrm{Log}}\).  The internal core therefore respects the branch
domain of \Cref{def:eml} at every step.
\end{proof}

\begin{proposition}[External public EML representation anchor]
\label{thm:eml-universality}
Odrzywo{\l}ek's public arXiv v2 record
\cite{Odrzywolek2026SingleBinaryOperator} asserts that every primitive in the
scientific-calculator basis \(\mathcal B_{36}\) displayed there has a finite
representative over \((\eml,1)\) on its stated principal-branch domain.  This
proposition is only the article's external public anchor for that assertion: the
thirty-six representatives and their domain checks are not printed or verified
inside this paper.  The anchor is used here only as an \(L0\)
representative-existence theorem; it supplies no normal form, no equality
algorithm for EML terms, no semantic compiler or normalizer, and no Richardson
witness identity except through the separate interface of \Cref{thm:separation}.
\end{proposition}

\begin{proof}
This is a citation step, not a term-by-term derivation.  The public anchor is the
versioned record arXiv:2603.21852v2, DOI
\texttt{10.48550/arXiv.2603.21852}.  The ancillary file named in \Cref{ass:HW}
is used only for the optional printed-string hypothesis; the main theorem chain
treats the full \(\mathcal B_{36}\) result as the single external existence
assertion stated here.  No later proof may appeal to an unprinted row of the
\(\mathcal B_{36}\) table; all self-contained EML calculations used later are
the core calculations isolated in \Cref{prop:internal-core-eml-witnesses}.
\end{proof}


\begin{theorem}[Exact scope of the \texorpdfstring{\(L0\)}{L0} import]
\label{thm:L0-citation-domain-support}
\label{thm:L0-external-anchor-criterion}
\label{thm:B36-import-exact-scope}
\label{thm:B36-import-normal-form-separation}
\label{cor:external-compiler-isolation}
\label{prop:B36-principal-branch-domain-source}
\label{thm:L0-import-replacement-invariance}
\label{thm:B36-citation-stability-and-domain-cut}
The imported Odrzywo{\l}ek input used in this paper is exactly the pinned
arXiv:2603.21852v2 claim that the main-text Table~1 scientific-calculator basis
\(\mathcal B_{36}\) has finite representatives over \((\eml,1)\) on the source's
stated principal-branch domains, with the versioned ancillary Part~II serving
only as the cited verification record for that representative-existence claim.
The ancillary Part~I Table~S2 strings identified in \Cref{ass:HW} are separate:
they enter this paper only through the optional printed-string hypothesis
\((H_{\mathcal W})\).  This pinning is part of the mathematical hypothesis:
replacing the citation by an unversioned web page, by another branch convention,
or by an endpoint-extended semantics would be a different assumption.  It
supplies representative existence only: it is not a term-by-term
self-contained proof in this article, not a branch-free or endpoint theorem, not
a canonical-form theorem, not an equality algorithm for EML terms, not a
semantic compiler, and not an endpoint-root Richardson package for \(\pi\), sine
on \(\RR\), and nonnegative square root on \(\RR_{\ge0}\).  Moreover, the later
Richardson, Zeckendorf \(L1\), one-orbit, and finite-cover \(L2\) arguments
remain valid if the full \(\mathcal B_{36}\) existence sentence is replaced by
the self-contained core witness package of
\Cref{prop:internal-core-eml-witnesses}; under that replacement the paper simply
no longer asserts \(L0\) expression universality for the whole calculator basis.
The multiple labels on this theorem are dependency-cut aliases for these exact
source, domain, normal-form, and replacement boundaries; they are not separate
imports from the cited table.
\end{theorem}

\begin{proof}
The citation \Cref{thm:eml-universality} imports one versioned public statement:
finite representatives for the main-text Table~1 \(\mathcal B_{36}\) basis, with
the principal logarithm and the domains fixed in that source.  The bibliographic
record pins the arXiv version, DOI, and ancillary file name; hence the manuscript
has a stable object to cite, but the object cited is still only the stated source
assertion and its cited Part~II verification record.  The manuscript does not
import any additional row, generated code artifact, or mutable repository state.
Every stronger item listed in the theorem would require extra data
absent from such an existence statement.  A branch-free assertion would compare
logarithm branches; an endpoint assertion would check boundary values; a
canonical form or equality algorithm would prove uniqueness or decidability for
EML terms; and a Richardson triple would prove the three displayed real-domain
identities.  None of these follows from mere representative existence, and none
is asserted by the citation step.

It remains to check that the full calculator-basis import is not hidden in the
strictness proofs.  The Richardson translation uses the core arithmetic library,
\(\exp\), division on its nonzero-denominator domain, positive square root, and a
internally repaired punctured witness package; these dependencies are recorded in
\Cref{lem:richardson-translation-bookkeeping,thm:separation}.  The Zeckendorf
protocol is proved from normal forms and value transport in
\Cref{thm:zeckendorf-protocol-univ}.  The one-orbit obstruction is the exponent
skeleton of \Cref{thm:level-zero-obstruction}.  The finite-cover statement uses
finite fibre data and the one-spike cover construction of
\Cref{thm:finite-state-strictness}.  Thus replacing the global calculator-basis
existence assertion by the internal core package preserves all strictness
arguments and removes only the headline claim that every operation in
\(\mathcal B_{36}\) has a cited representative.
\end{proof}

\begin{proposition}[Internal core EML witnesses]
\label{prop:internal-core-eml-witnesses}
In the principal-branch EML semantics, pure EML terms over \((\eml,1)\) represent
the constant \(1\), the identity input, \(\exp\), positive logarithm, negation,
addition, subtraction, multiplication, division by a nonzero real, rational
constants, \(\log2\), and square root on the positive domain \((0,\infty)\).
No value at \(0\) is asserted by this internal core package.
\end{proposition}

\begin{proof}
The explicit macros and real-domain calculations are printed in
\Cref{lem:eml-printed-core,lem:eml-real-witness-library}.  Those proofs use only
the definition of \(\eml\), positivity of the relevant inputs, and elementary
real algebra.
\end{proof}

\begin{theorem}[Self-contained \texorpdfstring{\(L0\)}{L0} core versus cited \texorpdfstring{\(\mathcal B_{36}\)}{B36} background]
\label{thm:self-contained-L0-core-cited-B36-boundary}
The self-contained \(L0\) content proved in this paper is exactly the finite core
witness package of \Cref{prop:internal-core-eml-witnesses}, together with the
calculated real-domain identities in
\Cref{lem:eml-printed-core,lem:eml-real-witness-library}.  The stronger statement
that every operation in Odrzywo{\l}ek's displayed calculator basis
\(\mathcal B_{36}\) has a finite representative over \((\eml,1)\) is used only as
the cited background theorem \Cref{thm:eml-universality}.  In particular, no
argument in the main hierarchy may invoke the \(\mathcal B_{36}\) import as a
self-contained term-by-term proof, as a normal-form algorithm, or as a verified
Richardson triple for \(\pi\), sine on \(\RR\), and square root on
\(\RR_{\ge0}\).
\end{theorem}

\begin{proof}
The first sentence is precisely the content of
\Cref{prop:internal-core-eml-witnesses}: its proof points to the printed macros
and real-domain calculations in
\Cref{lem:eml-printed-core,lem:eml-real-witness-library}, and no other
term-by-term EML representative calculation is performed in this section.  The
full calculator-basis assertion is stated separately in \Cref{thm:eml-universality},
whose proof imports the versioned public record rather than rederiving its
representatives.  \Cref{thm:B36-import-exact-scope} then proves that this import
has only representative-existence scope, and
\Cref{thm:L0-import-replacement-invariance} proves that replacing it by any
external theorem with the same existence content leaves the later \(L1\), \(L2\),
one-orbit, and punctured Richardson arguments unchanged.  The forbidden
uses listed in the last sentence are exactly the additional data excluded by
\Cref{thm:B36-import-exact-scope,cor:external-compiler-isolation,prop:B36-principal-branch-domain-source}:
term-by-term self-contained verification, canonical or decidable normal forms,
and the three real-domain identities required for the Richardson interface are
not consequences of the imported \(L0\) statement inside this paper.  Hence the
boundary is exact.
\end{proof}

\begin{theorem}[Submission replacement for a term-by-term \texorpdfstring{\(\mathcal B_{36}\)}{B36} table]
\label{thm:B36-self-contained-replacement-core}
For the theorem chain proved in this article, the only self-contained EML
representative data needed from \(L0\) are the finite core witnesses of
\Cref{prop:internal-core-eml-witnesses}.  Replacing the full imported
\(\mathcal B_{36}\) sentence of \Cref{thm:eml-universality} by the following
smaller internal assertion leaves every later theorem unchanged:
\[
\text{the core witnesses of \Cref{prop:internal-core-eml-witnesses} exist and
have the stated real-domain identities.}
\]
Thus a referee demand for self-containment is met inside the present proof graph
by shrinking the load-bearing \(L0\) library to this internally verified core; it
is not met, and is not claimed to be met, by an unprinted term-by-term proof of
all thirty-six \(\mathcal B_{36}\) representatives.
\end{theorem}

\begin{proof}
List the later dependencies on \(L0\).  The Richardson translation uses rational
constants, \(\log2\), arithmetic operations, \(\exp\), division on its nonzero
domain, and the positive square-root macro, together with a separately supplied
punctured triple for \(\pi\), sine, and absolute value off zero; these are exactly
the dependencies recorded in \Cref{lem:richardson-translation-bookkeeping} and
\Cref{thm:separation}.  The Zeckendorf protocol and semiring readouts in
\Cref{thm:zeckendorf-protocol-univ} are constructed from Fibonacci residues and
valuation transport, not from the calculator-basis representatives.  The
finite-fibre certificate and cover arguments use only finite sets, fibre counts,
Hankel matrices, Prony reconstruction, and the cover-size construction over
\(X_m\), as recorded in \Cref{prop:zeckendorf-cert-univ,thm:finite-state-strictness}.
Finally, the one-orbit obstruction is a statement about readouts from a free
monogenic orbit and prime exponent weights; it contains no EML representative
input.

The internal core representatives named in
\Cref{prop:internal-core-eml-witnesses} are proved by the printed calculations of
\Cref{lem:eml-printed-core,lem:eml-real-witness-library}.  Therefore substituting
that smaller assertion for the imported full \(\mathcal B_{36}\) existence
sentence preserves all hypotheses actually used by the Richardson interface,
the \(L1\) protocol, the \(L2\) cover strictness theorem, and the one-orbit
obstruction.  Conversely, \Cref{thm:B36-import-exact-scope} proves that the full
import supplies no normalizer, compiler, verified Richardson triple, or
term-by-term self-contained table inside this article.  Hence the replacement is
both sufficient for the submitted theorem chain and exact with respect to the
stronger \(\mathcal B_{36}\) self-containment claim.
\end{proof}

\begin{theorem}[Interface rigidity for the load-bearing \texorpdfstring{\(L0\)}{L0} and Richardson inputs]
\label{thm:L0-richardson-exact-load-bearing-cut}
\label{thm:submission-interface-rigidity}
The Proc. AMS theorem chain uses the expression-universality input and the
Richardson input through disjoint interfaces, and these interfaces are rigid in
the following sense.  Its \(L0\) interface is exactly one of the following two
choices:
\begin{enumerate}
\item the cited Odrzywo{\l}ek representative-existence theorem for the displayed
  \(\mathcal B_{36}\) basis, with no term-by-term verification reproduced here;
  or
\item the smaller self-contained core witness package of
  \Cref{prop:internal-core-eml-witnesses}, in which case the global
  \(\mathcal B_{36}\) expression-universality sentence is deleted but the later
  strictness theorems are unchanged.
\end{enumerate}
Its Richardson interface is exactly the punctured witness package of
\Cref{thm:separation}: finite pure EML witnesses for \(\pi\), sine on \(\RR\),
and \(|x|\) on \(\RR\setminus\{0\}\).  Consequently the cited
\(\mathcal B_{36}\) import, the finite syntax catalogue, and the endpoint-root
string \(U_{126}\) are not load-bearing substitutes for those identities.  In
particular, after deleting the optional printed-string hypothesis
\((H_{\mathcal W})\), every conclusion of \Cref{thm:main-hierarchy} that remains
in force factors through the replacement-triple Richardson interface; no
unconditional conclusion naming the printed pair \((U_{47},U_{111})\) remains
unless the two real-domain identities \(U_{47}=\pi\) and
\(U_{111}(x)=\sin x\) on \(\RR\) are separately supplied.
\end{theorem}

\begin{proof}
The two possible \(L0\) inputs are precisely the alternatives isolated in
\Cref{thm:L0-external-anchor-criterion}.  The first alternative is the cited
representative-existence proposition \Cref{thm:eml-universality}; by
\Cref{thm:B36-import-exact-scope} and \Cref{thm:B36-import-normal-form-separation}
it imports only existence of representatives on the stated principal-branch
domains, not a normalizer, equality procedure, compiler, or Richardson witness.
The second alternative is the replacement of the global calculator-basis claim by
the finite internal package of \Cref{prop:internal-core-eml-witnesses};
\Cref{thm:B36-self-contained-replacement-core} proves that every later
Richardson, Zeckendorf, one-orbit, and finite-cover argument still has all of its
actual hypotheses under that replacement.

It remains to identify the Richardson input.  In \Cref{thm:separation} and
\Cref{thm:replacement-triple-richardson-interface} the hypotheses are not the
existence of a large calculator-basis representation theorem.  They are the
finite punctured witnesses for \(\pi\), sine, and \(|x|\) off zero.  The proof
of \Cref{thm:B36-import-exact-scope} excludes deriving endpoint or normal-form
data from the \(\mathcal B_{36}\) import alone, and
\Cref{prop:printed-pi-sine-instantiation-criterion,cor:no-unconditional-printed-richardson-triple}
shows that the finite printed syntax of \((U_{47},U_{111})\) becomes a
Richardson instance exactly when those two printed strings have their
real-domain identities proved.  The Richardson theorem is relative to any
supplied punctured triple with those identities proved, and it proves the
absolute-value row internally, so the endpoint-root string \(U_{126}\) is not
load-bearing for \Cref{thm:main-hierarchy}.

Finally, \Cref{thm:main-hierarchy-exact-scope-closure} enumerates the four
interfaces through which the main theorem factors: the pinned \(L0\) import or
its internal-core replacement, the replacement-triple Richardson theorem, the
one-orbit \(L1\) obstruction, and the cover-relative \(L2\) budget obstruction.
The deletion of \((H_{\mathcal W})\) affects only the named printed
\((U_{47},U_{111})\) branch of the second interface.  Since all remaining
arrows in that enumeration cite the replacement-triple theorem rather than the
printed names, every surviving conclusion factors through the relative
interface, as claimed.
\end{proof}

\begin{theorem}[Load-bearing core after deleting optional external data]
\label{thm:load-bearing-core-dependency-cut}
Let the \emph{core theorem package} consist of the internally proved EML core of
\Cref{prop:internal-core-eml-witnesses}, the punctured Richardson interface
of \Cref{thm:separation}, the Zeckendorf \(L1\) protocol of
\Cref{thm:zeckendorf-protocol-univ}, the one-free-monogenic-orbit obstruction of
\Cref{thm:level-zero-obstruction}, the intrinsic \(L2\) transfer and height package of
\Cref{thm:intrinsic-fold-moment-transfer,thm:intrinsic-fold-fibre-height,cor:linfty-fold-entropy,thm:zero-temperature-fold-pressure},
and the cover-relative \(L2\) finite-state obstruction of
\Cref{thm:finite-state-strictness}.  If one deletes
\begin{enumerate}
\item the unprinted term-by-term \(\mathcal B_{36}\) representative table, keeping
  the internal core witnesses and a supplied punctured Richardson triple; and
\item the optional printed-string hypothesis \((H_{\mathcal W})\), keeping the
  verified-triple Richardson theorem for all triples whose \(\pi\), sine, and
  punctured absolute-value identities have been proved,
\end{enumerate}
then every assertion in the core theorem package remains proved with the same
statements.  The only assertions removed by this deletion are the global cited
\(\mathcal B_{36}\) expression-universality headline and the named printed
\((U_{47},U_{111})\) Richardson specialization.  In particular, the cover-relative
\(L2\) strictness theorem remains a theorem about computable covers over the fixed
normal-form sets \(X_m\), while the positive intrinsic \(\Fold_m\) fixed-\(q\)
transfer and height endpoint package remains the one proved directly in
\Cref{sec:L2}; no full intrinsic histogram or large-deviation classification of
the original fibres is introduced by the deletion.
\end{theorem}

\begin{proof}
The internal EML core remains because it is proved directly by
\Cref{lem:eml-printed-core,lem:eml-real-witness-library} and summarized in
\Cref{prop:internal-core-eml-witnesses}; it does not use the full
\(\mathcal B_{36}\) table.  The Richardson theorem used in the main spine is the
triple-relative \Cref{thm:separation}: its \(\pi\) and sine rows must be supplied
with proofs, and its punctured absolute-value row can be taken from the internal
term \(\Abs\).  Deleting \((H_{\mathcal W})\) therefore deletes only the optional
claim that the particular printed names \((U_{47},U_{111})\) supply the first two
rows.  It does not delete the universal verified-triple theorem, nor any instance
whose first two rows have been proved independently.  This is exactly the
interface identified in \Cref{thm:L0-richardson-exact-load-bearing-cut} and
\Cref{thm:unconditional-submission-richardson-closure}.

The Zeckendorf protocol theorem \Cref{thm:zeckendorf-protocol-univ} is proved from
finite normal forms, Fibonacci valuation transport, and Church-iteration
identities, not from EML calculator representatives or from the printed
Richardson strings.  The one-orbit obstruction \Cref{thm:level-zero-obstruction}
is a statement about readouts from one free monogenic orbit and its prime-weight
exponent skeleton, so it has no dependency on either deleted datum.  Finally,
\Cref{thm:intrinsic-fold-moment-transfer} proves the positive intrinsic transfer
statement directly from the bounded carry automaton.  Separately,
\Cref{thm:finite-state-strictness} is proved by the computable-cover skeleton:
\Cref{thm:finite-fiber-cover-realization} turns a computable positive fibre-size
function over \(X_m\) into a cover, \Cref{lem:cover-preserves-protocol} preserves
all normal-form readouts, and \Cref{prop:finite-state-coefficient-budget} gives
the exponential budget contradicted by a superexponential one-spike choice such
as \(a_m=m!\).  None of these cover-obstruction steps uses the intrinsic
multiplicities of \(\Fold_m\); by \Cref{thm:L2-cover-intrinsic-nonimplication} a
cover-moment transfer statement or obstruction cannot be transported to the
original fold fibres without an additional equality or comparison theorem.
Therefore the stated deletion leaves
exactly the listed core package and removes exactly the two optional headlines.
\end{proof}

\begin{definition}[Expression universality]\label{def:expr-univ}
A system is \(L0\)-expression-universal for a signature if every primitive
operation in the signature has at least one term representative over the chosen
primitive operations on its stated domain.  No normal form, equality procedure,
or fibre certificate is part of the definition.
\end{definition}

\begin{remark}\label{rem:eml-variants}
All EML claims in this article use the principal-branch convention of
\Cref{def:eml}; no branch-free complex calculus is asserted.
\end{remark}

\begin{remark}\label{rem:no-canon}
Expression universality alone does not choose canonical representatives.  This
is the gap separated by the Richardson obstruction in \Cref{thm:separation}.
\end{remark}


